\section*{Chapter \ref{ch:m2:mortgages-and-agencies} --- Mortgages and Agencies}

\begin{solution}{exo:m2:mortgages-and-agencies:1}
2\% ($10 \times 0.2\%$); 6\% ($1.5 \times 20 \times 0.2\%$); 12\% ($2 \times 6\%$,
the ramp being over after thirty months).
\end{solution}

\begin{solution}{exo:m2:mortgages-and-agencies:2}
$1 - 0.94^{1/12} = 0.5143\%$ a month; USD~257\,151 prepays.
\end{solution}

\begin{solution}{exo:m2:mortgages-and-agencies:3}
Agency, coupon, maturity, price, face value, settlement month. The seller
delivers the least valuable pools that meet the good-delivery guidelines, those
expected to prepay in the way that hurts the holder most. Pools worth more
than that, because their borrowers are expected to prepay more slowly, are kept
back and sold as \omterm{def:m2:mortgages-and-agencies:tba}{specified pools} at a pay-up.
\end{solution}

\begin{solution}{exo:m2:mortgages-and-agencies:4}
5.90\%. With no drop, rolling gives up the month's coupon (6\% a year on par)
and the paydown, returned at par rather than at 100.50; financing 100.50 of
cash at a little under the coupon is what that costs.
\end{solution}

\begin{solution}{exo:m2:mortgages-and-agencies:5}
At 4\% the incentive is two and a half points and the pool prepays at a CPR of
48\%: its principal comes back quickly at par and cannot be reinvested at 6\%.
The frozen-speed pool keeps paying 6\% on a balance that runs off at the
at-the-money speed, which is worth far more at a 4\% discount rate.
\end{solution}

\begin{solution}{exo:m2:mortgages-and-agencies:6}
At 5.5\%: 1.77 years and $-478$; at 6\%: 4.73 and $-560$; at 6.5\%: 6.51 and
$-93$. \omterm{def:m2:mortgages-and-agencies:convexity}{Negative convexity} is strongest near the money, where the S-curve is
steepest and a small rate move changes speeds the most; out of the money speeds
hardly respond and the pool behaves more like a bond.
\end{solution}

\begin{solution}{exo:m2:mortgages-and-agencies:7}
19.62, 14.73, 11.49, 7.73 and 5.77 years.
\end{solution}

\begin{solution}{exo:m2:mortgages-and-agencies:8}
The guarantee removes credit risk, not \omterm{def:m2:mortgages-and-agencies:cpr}{prepayment} risk. The holder is short
the borrowers' option to repay at par, so agency securities must yield more than
Treasuries of similar life to pay for that option, for their uncertain cash
flows and for their \omterm{def:m2:mortgages-and-agencies:convexity}{negative convexity}. For Fannie Mae and Freddie Mac the
guarantee is that of the enterprises, not the full faith and credit of the
United States that stands behind Ginnie Mae.
\end{solution}

\begin{solution}{pb:m2:mortgages-and-agencies:1}
\textbf{1.} CPR 11.9\%, price 100.00, \omterm{def:m2:mortgages-and-agencies:convexity}{effective duration} 4.73 years.
\textbf{2.} USD~4.73 million per basis point ($10^{10} \times 4.727 \times 10^{-4}$).
\textbf{3.} USD~73\,601.
\textbf{4.} About USD~6.42 billion.
\textbf{5.} $-560$: the duration shortens as rates fall and lengthens as they
rise, so the hedge must be traded in the direction of every move.
\textbf{6.} CPR 21.3\%, price 101.61, \omterm{def:m2:mortgages-and-agencies:convexity}{effective duration} 1.77 years.
\textbf{7.} USD~1.80 million per basis point.
\textbf{8.} By USD~2.93 million per basis point: the investor is short too much
duration and must buy it back.
\textbf{9.} About USD~3.97 billion, received fixed.
\textbf{10.} 6.51 years: the hedge would have to be increased instead.
\textbf{11.} About 1.4\%.
\textbf{12.} After a fall, hedgers receive fixed or buy bonds, pushing rates
lower still; after a rise, they pay or sell.
\textbf{13.} Investors with the opposite need, such as those who pay fixed
to hedge rising rates, and dealers, who charge for it; when they are scarce,
rates move further.
\textbf{14.} Options bought in advance deliver the duration automatically as
rates fall, so the investor trades less after the move; the convexity has been
passed to the option seller.
\textbf{15.} It bought pools and did not hedge them, so less of the market's
\omterm{def:m2:mortgages-and-agencies:convexity}{negative convexity} was being hedged dynamically.
\textbf{16.} The S-curve, its slope and the spread between mortgage and discount
rates are assumptions; real speeds also depend on loan age, size, borrowers'
credit and the lenders' capacity. The answer's order of magnitude is robust,
its digits are not.
\textbf{17.} The PSA-like ramp: new borrowers rarely refinance in their first
months whatever the incentive.
\textbf{18.} It loses fee income as the balance prepays, so its asset loses value
in a rally and it too must buy duration to hedge it.
\textbf{19.} Named result: \emph{the convexity hedge} for USD~10 billion after
a 50-basis-point rally is to receive fixed on about USD~4.0 billion of
ten-year swaps, as the \omterm{def:m2:mortgages-and-agencies:convexity}{effective duration} falls from 4.73 to 1.77 years.
\textbf{20.} Because borrowers refinance, their securities shorten, and hedged
holders must replace the duration they lost.
\end{solution}

\begin{solution}{iq:m2:mortgages-and-agencies:1}
Borrowers can repay at par at any time, so the holder is short a call option
on the loan. When rates fall the option moves into the money, \omterm{def:m2:mortgages-and-agencies:cpr}{prepayments}
accelerate and the price is capped near par: duration shortens. When rates
rise, \omterm{def:m2:mortgages-and-agencies:cpr}{prepayments} slow and duration lengthens. A price that rises less than it
falls is \omterm{def:m2:mortgages-and-agencies:convexity}{negative convexity}.
\iqlookfor{the embedded option and the duration changes it causes.}
\end{solution}

\begin{solution}{iq:m2:mortgages-and-agencies:2}
A forward trade in agency \omterm{def:m2:mortgages-and-agencies:mbs}{pass-throughs} on six parameters (agency, coupon,
maturity, price, face, settlement month), with the pools named two days before
settlement. It pools a million heterogeneous securities into a few fungible
contracts, which makes them liquid and cheap to trade; pools with valuable
characteristics trade separately as \omterm{def:m2:mortgages-and-agencies:tba}{specified pools}.
\iqlookfor{fungibility, and the \omterm{def:m2:bond-futures:ctd}{cheapest-to-deliver} consequence.}
\end{solution}

\begin{solution}{iq:m2:mortgages-and-agencies:3}
Sell the front-month TBA and buy the back month: the roll seller gives up the
coupon and paydown for a month and keeps the cash. The drop implies a financing
rate; when it is below repo, the roll is special. If you own the securities,
roll them and invest the cash; if you need the front-month pools, expect to pay
for them.
\iqlookfor{the implied financing rate and the comparison with repo.}
\end{solution}

\begin{solution}{iq:m2:mortgages-and-agencies:4}
Separate the drivers: house sales (seasoning, season of the year), refinancing
as a function of the incentive, burnout, loan size, credit, and lenders'
capacity; fit on loan-level or pool-level history; validate out of sample,
especially through past refinancing waves, and on the dispersion of speeds
across pools, not only their average. Watch that the model's durations hedge
well in practice.
\iqlookfor{the drivers, out-of-sample validation and a hedging test.}
\end{solution}

\begin{solution}{iq:m2:mortgages-and-agencies:5}
After a rally, mortgage durations shorten and hedgers receive fixed or buy
Treasuries; after a sell-off they pay or sell. The flows are in the direction
of the move, so they add momentum and raise realised and implied volatility,
particularly when much of the market is near the money. Central bank holdings
and options hedging reduce the effect.
\iqlookfor{direction, size, and what damps it.}
\end{solution}

\begin{solution}{iq:m2:mortgages-and-agencies:6}
Group pools by characteristics that drive speeds, or use the \omterm{def:m2:mortgages-and-agencies:cpr}{prepayment} model
vectorised over pools; reuse a common set of rate scenarios (up, down and base,
or simulated paths) and compute cash flows once per scenario; parallelise across
pools; cache results for pools whose inputs did not change; recompute only the
market-dependent parts intraday. Reconcile against a slower full run nightly,
and monitor the model's durations against realised hedge performance.
\iqlookfor{vectorisation, scenario reuse, caching and reconciliation.}
\end{solution}
